\section{Results}
\label{sec:results}

This section reports what was measured. Interpretation is deferred to
Section~\ref{sec:discussion}, and the wording is deliberately narrow: a difference is called
\emph{separated} only when its bootstrap interval excludes zero, and an ordering of means is
called an ordering of means.

\subsection{The protocol as it was run}
\label{sec:res-protocol}

Every rung is evaluated by pooled five-fold out-of-fold macro-F1. Each fold's held-out
predictions are collected, the five folds are concatenated into one prediction vector covering
every edge exactly once, and macro-F1 is computed on that vector. The metric covers all ten
classes on NF-UNSW-NB15 and eight of ten on NF-ToN-IoT, where \texttt{dos} (4 edges) and
\texttt{ransomware} (3 edges) remain in the graph for message passing but cannot survive a
five-way split. Because the class counts differ, the two datasets' absolute levels are not
comparable. Only the shape of the ladder is.

Each rung was trained at three seeds, 42, 1 and 2, with the fold partition held fixed at the
seed-42 stratified split. Only training stochasticity varies, so the pooled out-of-fold edge
set is identical across seeds, which is what makes the seed-matched bootstrap a paired
comparison. Reported standard deviations are training-seed variance at a fixed partition.
Fold-partition variance was not measured, so every $\pm$ in this section is a lower bound on
total uncertainty.

Two bootstrap procedures are used, both at 2{,}000 iterations. The primary is two-level: it
resamples edges and independently draws a training seed, so seed-to-seed variation enters the
interval. The secondary is seed-matched, pairing seed 42 with seed 42 and so on, resampling
edges only. The seed-matched intervals are narrower by construction and are reported so the
reader can see how much of an interval's width is seed variance. The decision rule was fixed
before the runs: separated if the 95\% two-level interval excludes zero, not separated
otherwise. Where the two procedures disagree, the two-level result decides and both are shown.

The semantic rung is deterministic given the folds and the frozen embeddings. Its standard
deviation is exactly 0.0000 at three seeds, and it enters every comparison as a fixed side.
All runs export \texttt{OMP\_NUM\_THREADS=1}. That pin is not cosmetic: unpinned runs drift by
roughly 0.012 macro-F1 at a fixed seed, which is larger than several of the differences
reported below, and one earlier finding in this project was retracted for exactly that reason.
Re-running the graph and semantic rungs at seed 42 reproduces their committed values exactly.
That has not always been true. A pre-v2 NF-ToN-IoT ladder failed to reproduce on two of four
rungs in August 2026, and thread pinning became mandatory afterwards.

Finally, one boundary carries through everything below. Edges are aggregated by source IP,
destination IP and attack class, so the label participates in defining edge identity. These are
controlled comparisons between representations and mechanisms, not deployment-valid detection
estimates.

\subsection{The model ladder}
\label{sec:res-ladder}

Table~\ref{tab:ladder} gives the five rungs on both datasets. Alongside the four rungs of the
designed ladder sits \emph{head alone}: the per-fold classification head trained on the semantic
embeddings, evaluated on its own. It is the strongest single-modality baseline available, and it
is also the consultant the loop queries, so it is reported in every table rather than only in the
ablations.

\begin{table}[htbp]
\centering
\caption{Pooled out-of-fold macro-F1 by rung, three training seeds at a fixed fold partition.
The $\pm$ is training-seed variance only. Absolute levels are not comparable across datasets.}
\label{tab:ladder}
\small
\resizebox{\textwidth}{!}{%
\begin{tabular}{lcccc}
\toprule
& \multicolumn{2}{c}{NF-UNSW-NB15 (10 classes)} & \multicolumn{2}{c}{NF-ToN-IoT (8 classes)} \\
\cmidrule(lr){2-3}\cmidrule(lr){4-5}
Rung & mean $\pm$ std & seeds 42 / 1 / 2 & mean $\pm$ std & seeds 42 / 1 / 2 \\
\midrule
Graph encoder      & $0.7437 \pm 0.0228$ & 0.7219 / 0.7674 / 0.7418 & $0.4336 \pm 0.0053$ & 0.4290 / 0.4393 / 0.4326 \\
Semantic rung      & $0.7353 \pm 0.0000$ & 0.7353 / 0.7353 / 0.7353 & $0.2785 \pm 0.0000$ & 0.2785 / 0.2785 / 0.2785 \\
AGAF               & $0.7793 \pm 0.0116$ & 0.7680 / 0.7912 / 0.7788 & $0.4102 \pm 0.0279$ & 0.3975 / 0.3910 / 0.4422 \\
Feedback loop      & $0.8341 \pm 0.0042$ & 0.8332 / 0.8387 / 0.8304 & $0.5044 \pm 0.0080$ & 0.4985 / 0.5012 / 0.5135 \\
Head alone         & $0.8374 \pm 0.0048$ & 0.8321 / 0.8386 / 0.8416 & $0.5081 \pm 0.0073$ & 0.5165 / 0.5035 / 0.5043 \\
\bottomrule
\end{tabular}
}
\end{table}

By mean, the order is identical on the two datasets at the top and differs in the middle. On
NF-UNSW-NB15 it runs head alone, loop, AGAF, graph encoder, semantic rung. On NF-ToN-IoT it runs
head alone, loop, graph encoder, AGAF, semantic rung: AGAF falls below the rung it is supposed to
improve on. The semantic rung is the weakest on both, and the gap is far larger under NF-ToN-IoT's
imbalance, where the embedding-only prototype scorer reaches 0.2785 against the graph encoder's
0.4336.

Table~\ref{tab:comparisons} answers RQ1 and RQ2 by turning those orderings into intervals.

\begin{table}[htbp]
\centering
\caption{Rung comparisons, both bootstrap procedures, 2{,}000 iterations. \emph{Sep} marks an
interval that excludes zero. \emph{Stable} reports whether the sign of the per-seed difference
is the same at all three seeds.}
\label{tab:comparisons}
\footnotesize
\resizebox{\textwidth}{!}{%
\begin{tabular}{llcccc}
\toprule
Dataset & Comparison & two-level mean [95\% CI] & sep & seed-matched mean [95\% CI] & stable \\
\midrule
\multirow{6}{*}{UNSW}
 & loop $-$ AGAF        & $+0.0546$ [$+0.0228$, $+0.0880$] & yes & $+0.0551$ [$+0.0331$, $+0.0779$] & yes \\
 & loop $-$ graph       & $+0.0908$ [$+0.0412$, $+0.1407$] & yes & $+0.0907$ [$+0.0571$, $+0.1250$] & yes \\
 & loop $-$ semantic    & $+0.0989$ [$+0.0721$, $+0.1267$] & yes & $+0.0991$ [$+0.0743$, $+0.1253$] & yes \\
 & head $-$ graph       & $+0.0938$ [$+0.0420$, $+0.1413$] & yes & $+0.0939$ [$+0.0607$, $+0.1278$] & yes \\
 & AGAF $-$ graph       & $+0.0362$ [$-0.0096$, $+0.0822$] & no  & $+0.0356$ [$+0.0038$, $+0.0678$] & yes \\
 & loop $-$ head        & $-0.0030$ [$-0.0242$, $+0.0136$] & no  & $-0.0032$ [$-0.0121$, $+0.0064$] & no \\
\midrule
\multirow{6}{*}{ToN}
 & loop $-$ AGAF        & $+0.0932$ [$+0.0180$, $+0.1667$] & yes & $+0.0940$ [$+0.0406$, $+0.1480$] & yes \\
 & loop $-$ graph       & $+0.0694$ [$+0.0129$, $+0.1288$] & yes & $+0.0691$ [$+0.0171$, $+0.1200$] & yes \\
 & loop $-$ semantic    & $+0.2231$ [$+0.1703$, $+0.2771$] & yes & $+0.2228$ [$+0.1732$, $+0.2713$] & yes \\
 & head $-$ graph       & $+0.0723$ [$+0.0125$, $+0.1322$] & yes & $+0.0725$ [$+0.0224$, $+0.1236$] & yes \\
 & AGAF $-$ graph       & $-0.0238$ [$-0.0921$, $+0.0607$] & no  & $-0.0249$ [$-0.0697$, $+0.0226$] & no \\
 & loop $-$ head        & $-0.0028$ [$-0.0388$, $+0.0318$] & no  & $-0.0034$ [$-0.0224$, $+0.0129$] & no \\
\bottomrule
\end{tabular}
}
\end{table}

Four results hold on both datasets. The loop is separated above AGAF, above the graph encoder
and above the semantic rung, and the head alone is separated above the graph encoder. This is
the first ladder step in this project's history that survives training-seed variance, and it
survives on two datasets with the same architecture and the same folds.

Two results are negative, also on both datasets. AGAF is not separated from the graph encoder:
positive on NF-UNSW-NB15 and negative on NF-ToN-IoT, with an interval crossing zero either way.
The seed-matched interval for NF-UNSW-NB15 does exclude zero, $+0.0356$ [$+0.0038$, $+0.0678$],
which is the clearest illustration in this report of what the two-level procedure is for: once
seed variance enters, the same difference stops being separated. Under the pre-registered rule,
AGAF versus the graph encoder is not separated, and the answer to RQ2 is negative on both
datasets.

The second negative result is the one that matters most. The loop is not separated from the head
it consults. The difference is $-0.0030$ [$-0.0242$, $+0.0136$] on NF-UNSW-NB15 and $-0.0028$
[$-0.0388$, $+0.0318$] on NF-ToN-IoT. The head carries the higher mean on both, the sign is
unstable across seeds, and at NF-UNSW-NB15 seed 42 the ordering flips, with the loop at 0.8332
against the head's 0.8321. The right word is indistinguishable. No statement of the loop's
margin over AGAF or the graph encoder in this report appears without this sentence beside it.

\subsection{Feedback ablations and the mechanism in isolation}
\label{sec:res-ablations}

RQ3 asks whether the feedback stage adds anything beyond static fusion, and if so, through which
path. Two experiments answer it, and they answer differently.

The first is the ladder-level ablation, in which the full system runs with real advice, with the
advice channel switched off (\texttt{head\_only}), and with random advice. Fusion at the output
is active in all three arms. Table~\ref{tab:ablations} gives every seed.

\begin{table}[htbp]
\centering
\caption{Feedback ablation, per training seed, fusion active in all arms. The interval is a
paired bootstrap of real against \texttt{head\_only}.}
\label{tab:ablations}
\small
\resizebox{\textwidth}{!}{%
\begin{tabular}{llcccc}
\toprule
Dataset & Seed & real & \texttt{head\_only} & random & real $-$ \texttt{head\_only} [95\% CI] \\
\midrule
\multirow{3}{*}{UNSW}
 & 42 & 0.8332 & 0.7543 & 0.6385 & $+0.0792$ [$+0.0409$, $+0.1185$] \\
 & 1  & 0.8387 & 0.7336 & 0.6511 & $+0.1054$ [$+0.0671$, $+0.1437$] \\
 & 2  & 0.8304 & 0.7314 & 0.6378 & $+0.0994$ [$+0.0595$, $+0.1387$] \\
\midrule
\multirow{3}{*}{ToN}
 & 42 & 0.4985 & 0.3809 & 0.4002 & $+0.1159$ [$+0.0596$, $+0.1751$] \\
 & 1  & 0.5012 & 0.4404 & 0.3790 & $+0.0595$ [$+0.0074$, $+0.1148$] \\
 & 2  & 0.5135 & 0.4374 & 0.4067 & $+0.0743$ [$+0.0223$, $+0.1280$] \\
\bottomrule
\end{tabular}
}
\end{table}

Feedback on scores above feedback off at every seed on both datasets, and every interval excludes
zero. Random advice costs at least 0.09 macro-F1 against real advice at every seed, so the arm is
sensitive to what the advice says and not merely to the presence of an extra tensor. What this
comparison cannot do is attribute the gain, because the real arm differs from
\texttt{head\_only} in two ways at once: semantic evidence re-enters the edge representations,
and semantic logits are fused at the output.

The second experiment separates them. Output fusion is disabled in every arm, so advice can
only reach the classifier through the injection path under test. Table~\ref{tab:mechanism}
gives the result.

\begin{table}[htbp]
\centering
\caption{Mechanism in isolation, output fusion disabled in every arm. Means are over three
seeds; intervals are paired over the fifteen (seed, fold) pairs. Captured share is the arm's
gain as a fraction of the oracle's.}
\label{tab:mechanism}
\footnotesize
\resizebox{\textwidth}{!}{%
\begin{tabular}{lcccc}
\toprule
& \multicolumn{2}{c}{NF-UNSW-NB15} & \multicolumn{2}{c}{NF-ToN-IoT} \\
\cmidrule(lr){2-3}\cmidrule(lr){4-5}
Arm & mean & paired diff vs control [95\% CI] & mean & paired diff vs control [95\% CI] \\
\midrule
control (\texttt{head\_only}) & 0.7398 & -- & 0.4195 & -- \\
prototype, edge   & 0.7382 & $-0.0046$ [$-0.0273$, $+0.0158$] & 0.4325 & $+0.0099$ [$-0.0111$, $+0.0286$] \\
trained head, edge & 0.7441 & $+0.0007$ [$-0.0183$, $+0.0201$] & 0.3781 & $-0.0378$ [$-0.0588$, $-0.0167$] \\
oracle, edge      & 0.7743 & $+0.0369$ [$+0.0097$, $+0.0618$] & 0.5286 & $+0.1354$ [$+0.0919$, $+0.1797$] \\
oracle, attention & 0.7382 & $-0.0007$ [$-0.0211$, $+0.0180$] & 0.4310 & $+0.0114$ [$-0.0083$, $+0.0307$] \\
\midrule
\multicolumn{5}{l}{\emph{Pooled headroom and captured share}} \\
oracle headroom   & \multicolumn{2}{c}{$+0.0345$} & \multicolumn{2}{c}{$+0.1090$} \\
prototype share   & \multicolumn{2}{c}{$-4.5\%$} & \multicolumn{2}{c}{$11.9\%$} \\
trained head share & \multicolumn{2}{c}{$12.5\%$} & \multicolumn{2}{c}{$-38.0\%$} \\
\bottomrule
\end{tabular}
}
\end{table}

The oracle receives the true label at $\pm 4$ nats through the same projection every other arm
uses. It is a deliberate leak and is a diagnostic only, never a performance figure. It gains
$+0.0369$ [$+0.0097$, $+0.0618$] on NF-UNSW-NB15 and $+0.1354$ [$+0.0919$, $+0.1797$] on
NF-ToN-IoT, both separated. The channel therefore has measurable capacity on both graphs.

Neither realistic consultant uses it. The embedding-only prototype captures $-4.5\%$ of the
NF-UNSW-NB15 headroom and $11.9\%$ of the NF-ToN-IoT headroom, neither separated from zero. The
trained head, which scores 0.8374 and 0.5081 as a standalone classifier, captures $12.5\%$ on
NF-UNSW-NB15, again not separated, and $-38.0\%$ on NF-ToN-IoT, which is a separated regression.
A cross-fitted variant of the same head, measured at seed 42 only, moves the arm from $-0.0230$
to $-0.0654$ against control on NF-UNSW-NB15 and from $-0.0165$ to $+0.0070$ on NF-ToN-IoT, so
matching the advice's train-time and test-time reliability does not rescue it either.

One caveat travels with the NF-UNSW-NB15 headroom. At seed 42 alone the oracle arm scores
0.7522 against a control of 0.7543, a headroom of $-0.0021$, so a captured share is not
computable there. The $+0.0345$ is a three-seed mean, and single-seed shares on NF-UNSW-NB15 are
not reported.

Read together, the two experiments give RQ3 a split answer. Feedback plus fusion is separated
from fusion alone at every seed on both datasets. The injection mechanism alone, driven by any
consultant we can actually build, is not separated from doing nothing.

\subsection{Five attempts on the embedding-only consultant}
\label{sec:res-prototype}

Before the consultant was changed, five treatments were run against the whitened-prototype
scorer. Each attacked a different candidate explanation for why its advice does not convert.
All five failed, and the pattern of failure is the evidence that the constraint is the
consultant rather than the mechanism around it.

Two of them were signal-quality fixes, measured at three seeds as conditions A, B and C.
Condition A is the canonical configuration. Condition B trains the selector head whose entropy
ranks edges, and calibrates the consultant's temperature per fold instead of pinning it. Condition
C re-selects the knobs on top of B. The two defects were genuinely removed: on NF-UNSW-NB15 the
selector head's macro-F1 moved from 0.019--0.036 to 0.646--0.684, mean disagreement fell from
0.8975 to 0.524--0.548, and the consultant's maximum softmax probability rose from 0.102 to
0.487, with the temperature falling from about 9.4 to 0.105. The loop then scored 0.7644 in A,
0.7610 in B and 0.7809 in C on NF-UNSW-NB15, and 0.4521, 0.4147 and 0.4630 on NF-ToN-IoT.
Fixing both signals made the loop worse on both datasets. Condition C recovers, but its
selection curves are flat, NF-UNSW-NB15's argmax landed on the boundary of the swept range at
$k = 35$, and \texttt{loop\_vs\_agaf} includes zero in all three conditions.

The same runs produced the clearest single diagnostic in the feedback work. The magnitude
actually added to the 39-wide encoded edge vector, whose features have unit variance, was
0.077 / 0.060 / 0.061 on NF-UNSW-NB15 and 0.023 / 0.021 / 0.023 on NF-ToN-IoT under condition A.
Calibrating the temperature raised it to roughly 1.2 and 2.8 respectively, a factor of about 17
and about 125. The canonical prototype loop was, in effect, injecting almost nothing.

Table~\ref{tab:treatments} lists all five treatments and what each was worth.

\begin{table}[htbp]
\centering
\caption{Five treatments applied to the embedding-only consultant. None was kept. Rows marked
seed 42 carry no interval across seeds and are single-seed orderings only.}
\label{tab:treatments}
\footnotesize
\begin{tabularx}{\textwidth}{@{}l X X c@{}}
\toprule
Treatment & What it changed & Result & Kept \\
\midrule
Selector-head supervision & the ranking head is trained as a classifier & 3 seeds: loop worse on both datasets & no \\
Consultant temperature & per-fold calibration instead of a pinned $T$ & 3 seeds: loop worse, jointly with the above & no \\
Reliability weighting & gate and magnitude scaled by per-class train precision & seed 42: $+0.0057$ / $+0.0094$, below the 0.02 rule & no \\
Advice format & one-hot and softmax at the oracle's $\pm 4$ nats & seed 42: every arm below baseline on both & no \\
Cross-fitted advice & train rows replaced by inner out-of-fold logits & seed 42: $-0.0246$ / $-0.0338$ against in-fold & no \\
\bottomrule
\end{tabularx}
\end{table}

The advice-format arm produced a counter-intuitive measurement worth stating on its own. Handing
the consultant the oracle's saturated format made the input magnitude larger and the injected
magnitude smaller, from 0.0767 to 0.0183 on NF-UNSW-NB15 and from 0.0234 to 0.0202 on
NF-ToN-IoT. The projection is learned, so it adapts to a larger and lower-variance input by
injecting less. Whatever the oracle's advantage is, it is not the shape of the vector it presents.

The cross-fitting arm tested the loop's trust directly. If in-fold optimism in the advice were
what the loop was exploiting, honest advice should lower the scalar the loop learns for how hard
to push. Across four settings the learned trust scalar moved by at most 0.0205, from 1.5722 to
1.5664 and 1.5227 to 1.5022 on NF-UNSW-NB15 and from 1.4768 to 1.4652 and 1.4842 to 1.4899 on
NF-ToN-IoT, moving the wrong way in the last of those. What did move was the injected magnitude,
from 1.3430 to 1.1752 and from 1.3291 to 1.0745. The model absorbs the change in the projection
rather than in the parameter that is supposed to represent trust.

Finally, the consultants differ in how well they complement the graph encoder on the edges the
loop actually consults. On the flagged set at seed 42 the graph encoder is 0.603 accurate on
NF-UNSW-NB15 and 0.359 on NF-ToN-IoT. The prototype scorer reaches 0.672 and 0.348 there, so on
NF-ToN-IoT it is no better than the model it is meant to correct. The trained head reaches 0.770
and 0.759, and the resulting loop reaches 0.750 and 0.724. On NF-ToN-IoT, among the flagged edges
where the head and the graph encoder disagree, the head is right 245 times and wrong 32, and
inside the confidence-gated half the split is 145 to 4.
\emph{[GAP: the same disagreement split for NF-UNSW-NB15, and the equivalent for the prototype
consultant on either dataset, was never recorded as an artifact.]}

\subsection{Attention injection}
\label{sec:res-attention}

The alternative injection design biases graph attention rather than the edge representation. It
produces churn of exactly 0.0000 in every configuration tested, including when handed the true
label at $\pm 4$ nats. Not one prediction changes. Table~\ref{tab:mechanism} confirms it a third
time under the current encoding: the oracle attention arm is $-0.0007$ [$-0.0211$, $+0.0180$] on
NF-UNSW-NB15 and $+0.0114$ [$-0.0083$, $+0.0307$] on NF-ToN-IoT, neither separated from zero.

The cause is structural and holds independently of the encoding. Attention modifies node
embeddings, but 58.7\% of NF-UNSW-NB15 edges share a source and destination pair with at least
one other edge, so node-derived terms are identical for them and only the edge attributes
distinguish them, which attention never touches. On NF-ToN-IoT a different mechanism produces the
same outcome: 62.8\% of edges have a destination of in-degree one, where a softmax over incoming
edges is 1.0 whatever bias is applied. Two datasets, two independent reasons, the same null.

\subsection{Re-trained published baselines}
\label{sec:res-baselines}

RQ5 compares the system against two published graph-based intrusion detection architectures,
E-GraphSAGE and TE-G-SAGE, re-trained on the same aggregated rows, folds, evaluation classes and
metric. These are not comparisons against the papers' published numbers, which were obtained on
per-flow graphs roughly 300 times larger. Table~\ref{tab:baselines} gives the point estimates.

\begin{table}[htbp]
\centering
\caption{Re-trained baselines, mean pooled out-of-fold macro-F1 over seeds 42, 1 and 2.
\emph{Faithful} marks a configuration that is the published architecture; \texttt{+ node
features} is a labelled ablation of ours and must not be tabulated under the paper's name.}
\label{tab:baselines}
\small
\resizebox{\textwidth}{!}{%
\begin{tabular}{lllcc}
\toprule
Architecture & Configuration & Faithful & NF-UNSW-NB15 & NF-ToN-IoT \\
\midrule
\multirow{3}{*}{E-GraphSAGE}
 & as published        & yes & $0.1194 \pm 0.0065$ & $0.4141 \pm 0.0046$ \\
 & refit               & yes & $0.1368 \pm 0.0076$ & $0.4141 \pm 0.0046$ \\
 & $+$ our node features & no & $0.1173 \pm 0.0065$ & $0.4358 \pm 0.0323$ \\
\midrule
\multirow{3}{*}{TE-G-SAGE}
 & as published        & yes & $0.3841 \pm 0.0078$ & $0.1931 \pm 0.0161$ \\
 & refit               & yes & $0.5842 \pm 0.0224$ & $0.2345 \pm 0.0101$ \\
 & $+$ our node features & no & $0.7196 \pm 0.0094$ & $0.3523 \pm 0.0182$ \\
\bottomrule
\end{tabular}
}
\end{table}

The collapse of E-GraphSAGE on NF-UNSW-NB15 is a representation effect with a testable
prediction attached, not evidence that the architecture is weak. Its edge classifier reads only the
concatenation $[h_u \Vert h_v]$ of the two endpoint embeddings. An edge's own attributes never
reach it, so parallel edges between one address pair cannot be told apart. That predicts a
collapse
proportional to how common parallel edges are. On NF-UNSW-NB15, 385 of 656 edges are
indistinguishable by endpoints, that is 58.7\%, and the architecture scores 0.1194. On
NF-ToN-IoT only 167 of 2{,}127 edges are, that is 7.9\%, and the same architecture scores
0.4141. The prediction and the inversion agree.
The ceiling is also not a training-schedule artifact: trained to the authors' full 4{,}999
epochs on fold 0, validation macro-F1 plateaus near 0.11 and peaks at 0.1384 at epoch 1{,}400.

Table~\ref{tab:vsbaselines} gives the intervals for the two strongest rungs.

\begin{table}[htbp]
\centering
\caption{Loop and head alone against each re-trained baseline, two-level bootstrap over edges
and both seed sets. Positive favours our rung. \emph{ns} marks an interval containing zero.}
\label{tab:vsbaselines}
\footnotesize
\resizebox{\textwidth}{!}{%
\begin{tabular}{llll}
\toprule
Dataset & Baseline & Loop $-$ baseline & Head alone $-$ baseline \\
\midrule
\multirow{6}{*}{UNSW}
 & E-GraphSAGE, as published & $+0.7130$ [$+0.6708$, $+0.7540$] & $+0.7161$ [$+0.6744$, $+0.7583$] \\
 & E-GraphSAGE, refit        & $+0.6960$ [$+0.6526$, $+0.7373$] & $+0.6991$ [$+0.6545$, $+0.7403$] \\
 & E-GraphSAGE $+$ features  & $+0.7157$ [$+0.6731$, $+0.7558$] & $+0.7188$ [$+0.6759$, $+0.7598$] \\
 & TE-G-SAGE, as published   & $+0.4500$ [$+0.3924$, $+0.5042$] & $+0.4531$ [$+0.3976$, $+0.5052$] \\
 & TE-G-SAGE, refit          & $+0.2519$ [$+0.1914$, $+0.3169$] & $+0.2550$ [$+0.1943$, $+0.3197$] \\
 & TE-G-SAGE $+$ features    & $+0.1155$ [$+0.0673$, $+0.1627$] & $+0.1186$ [$+0.0691$, $+0.1659$] \\
\midrule
\multirow{6}{*}{ToN}
 & E-GraphSAGE, as published & $+0.0890$ [$+0.0312$, $+0.1480$] & $+0.0925$ [$+0.0361$, $+0.1506$] \\
 & E-GraphSAGE, refit        & $+0.0890$ [$+0.0312$, $+0.1480$] & $+0.0925$ [$+0.0361$, $+0.1506$] \\
 & E-GraphSAGE $+$ features  & $+0.0667$ [$-0.0181$, $+0.1505$] \emph{ns} & $+0.0702$ [$-0.0139$, $+0.1512$] \emph{ns} \\
 & TE-G-SAGE, as published   & $+0.3083$ [$+0.2345$, $+0.3793$] & $+0.3119$ [$+0.2361$, $+0.3834$] \\
 & TE-G-SAGE, refit          & $+0.2672$ [$+0.2015$, $+0.3348$] & $+0.2708$ [$+0.2034$, $+0.3362$] \\
 & TE-G-SAGE $+$ features    & $+0.1502$ [$+0.0749$, $+0.2268$] & $+0.1537$ [$+0.0796$, $+0.2246$] \\
\bottomrule
\end{tabular}
}
\end{table}

The loop and the head alone are separated above every faithful re-trained baseline on both
datasets. The only cells that are not separated involve \texttt{+ our node features}, which is
our ablation and not a published architecture, and even those are separated under the
seed-matched procedure. The lower rungs are weaker: on NF-ToN-IoT neither the graph encoder
($+0.0203$ [$-0.0230$, $+0.0633$]) nor AGAF ($-0.0061$ [$-0.0797$, $+0.0777$]) is separated from
E-GraphSAGE as published, and on NF-UNSW-NB15 neither the graph encoder ($+0.0249$) nor the
semantic rung ($+0.0163$) is separated from TE-G-SAGE with our node features. Note also that the
loop's margin over any baseline stays within about 0.004 of the head alone's, everywhere in the
table. That is the unseparated \mbox{loop--head} difference seen from the baseline side.

Because most of the NF-UNSW-NB15 margin is not architectural, it is worth decomposing.
Table~\ref{tab:decomposition} walks from the strongest published configuration to our rungs.
Every level is a three-seed mean, no step carries an interval, and no step is a separation claim.

\begin{table}[htbp]
\centering
\caption{Decomposition of the NF-UNSW-NB15 gap against TE-G-SAGE. Point estimates only.
Shares are of the total gap to the loop.}
\label{tab:decomposition}
\small
\resizebox{\textwidth}{!}{%
\begin{tabular}{llcc}
\toprule
Step & macro-F1 & attributable to & share \\
\midrule
TE-G-SAGE as published        & 0.3841 & -- & -- \\
$+$ fair tuning (refit)       & 0.5842 & $+0.2001$, graph scale and schedule & 44.5\% \\
$+$ our ten centralities      & 0.7196 & $+0.1354$, feature engineering      & 30.1\% \\
our loop rung                 & 0.8341 & $+0.1145$, architecture            & 25.4\% \\
\midrule
(our AGAF rung)               & 0.7793 & $+0.0597$ over the same base       & 15.1\% of 0.3952 \\
\bottomrule
\end{tabular}
}
\end{table}

Roughly 45\% of the naive gap is our choice of how to tune, roughly 30\% is our node features,
and roughly 25\% is the architecture. Measured to AGAF instead of the loop, the architecture
share is 15.1\%.

\subsection{Per-class behaviour}
\label{sec:res-perclass}

Macro-F1 averages over classes of very different sizes, so Table~\ref{tab:perclass} gives the
per-class picture with the class sizes beside it.

\begin{table}[htbp]
\centering
\caption{Per-class F1, mean over three seeds, with edge counts. NF-ToN-IoT's \texttt{dos} and
\texttt{ransomware} are excluded from the metric and omitted here.}
\label{tab:perclass}
\small
\resizebox{\textwidth}{!}{%
\begin{tabular}{lrcccc}
\toprule
Class & edges & graph & AGAF & loop & head alone \\
\midrule
\multicolumn{6}{l}{\emph{NF-UNSW-NB15}} \\
Normal          & 311 & 0.9129 & 0.9152 & 0.9543 & 0.9617 \\
Analysis        & 25  & 0.4327 & 0.5886 & 0.6196 & 0.6407 \\
Backdoors       & 40  & 0.5929 & 0.6578 & 0.8019 & 0.7983 \\
DoS             & 40  & 0.5831 & 0.6446 & 0.7424 & 0.7311 \\
Exploits        & 40  & 0.9309 & 0.9583 & 0.9580 & 0.9500 \\
Fuzzers         & 40  & 0.6207 & 0.7330 & 0.7749 & 0.7939 \\
Generic         & 40  & 0.6591 & 0.5409 & 0.6899 & 0.6955 \\
Reconnaissance  & 40  & 0.9110 & 0.8916 & 0.8731 & 0.8865 \\
Shellcode       & 40  & 0.8669 & 0.9428 & 0.9959 & 1.0000 \\
Worms           & 40  & 0.9272 & 0.9205 & 0.9311 & 0.9165 \\
\midrule
\multicolumn{6}{l}{\emph{NF-ToN-IoT}} \\
Benign          & 1746 & 0.9599 & 0.9275 & 0.9832 & 0.9840 \\
backdoor        & 16   & 0.3633 & 0.2439 & 0.4467 & 0.5064 \\
ddos            & 35   & 0.1848 & 0.2419 & 0.2913 & 0.3145 \\
injection       & 116  & 0.5436 & 0.4845 & 0.6277 & 0.5888 \\
mitm            & 157  & 0.8838 & 0.7082 & 0.8566 & 0.8339 \\
password        & 25   & 0.0401 & 0.2950 & 0.2100 & 0.2193 \\
scanning        & 13   & 0.0958 & 0.2781 & 0.3116 & 0.3004 \\
xss             & 12   & 0.3728 & 0.0844 & 0.2612 & 0.2678 \\
\bottomrule
\end{tabular}
}
\end{table}

The loop is above AGAF on eight of ten NF-UNSW-NB15 classes. The two exceptions are
Reconnaissance, 0.8731 against 0.8916, and Exploits, where the two are level to three decimals.
On NF-ToN-IoT the loop is above AGAF on seven of eight, losing only \texttt{password} at 0.2100
against 0.2950. Against the head alone the picture is a near-tie, as the pooled numbers imply.
The loop is above the head on four NF-UNSW-NB15 classes, Backdoors, DoS, Exploits and Worms, and
on three NF-ToN-IoT classes, \texttt{injection}, \texttt{mitm} and \texttt{scanning}. Every
NF-UNSW-NB15 difference between the two is smaller than 0.022. The largest single deficit against
the head anywhere is NF-ToN-IoT \texttt{backdoor}, 0.4467 against 0.5064, on a class of 16
edges.

Two rare-class observations matter for RQ4. AGAF's NF-ToN-IoT collapse is concentrated in the
classes where the graph encoder is strong: \texttt{mitm} falls from 0.8838 to 0.7082 and
\texttt{xss} from 0.3728 to 0.0844. Meanwhile AGAF is the only rung that lifts \texttt{password}
above 0.29 from the graph encoder's 0.0401. On the smallest classes, single edges move F1 by
large amounts, and none of these per-class differences carries an interval.

\subsection{Knob selection}
\label{sec:res-knobs}

Two quantities are selected per consultant and per dataset on validation folds only, never on
test: the entropy percentile that flags edges and the injection scale. The selected values are
$k = 29$ and scale 2.0 on NF-UNSW-NB15 and $k = 16$ and scale 20.0 on NF-ToN-IoT. With the
confidence gate keeping the top half of the flagged set, the effective consultation rate is
14.5\% of edges on NF-UNSW-NB15 and 8.0\% on NF-ToN-IoT. At seed 42 that is 204 of 656 edges
flagged and 102 gated on NF-UNSW-NB15, and 532 of 2{,}127 flagged and 266 gated on NF-ToN-IoT.

Neither quantity is tuned, and the report does not describe them that way. Averaged over the
three selection seeds, the 21 candidates for $k$ span 0.8238 to 0.8277 of validation macro-F1 on
NF-UNSW-NB15 and 0.5170 to 0.5232 on NF-ToN-IoT. The six injection-scale candidates span 0.8252
to 0.8277 and 0.5196 to 0.5232. Every one of those spans is smaller than the rungs' own
seed-to-seed spread, so the argmax is not distinguishable from its neighbours and the selection
is picking noise. NF-ToN-IoT's scale selected on the upper boundary of the swept range at 20.0,
which means the sweep cannot say whether an optimum lies outside it. The contract records that
boundary condition explicitly.
